\chapter{Oligopoly and Collusion}\label{ch:b3:oligopoly-and-collusion}

Most markets are served by a few firms selling products that differ: a handful of
brands of motor oil, three bakeries in a town, two airlines on a route. Their prices
depend on how close their products are, on how many firms the market can support and
on whether they compete or coordinate. This chapter builds the models used to answer
those questions with data --- differentiated demand, \omterm{def:b3:oligopoly-and-collusion:entry-threshold}{entry thresholds}, markup
estimators, collusion conditions --- and holds them against the record of mergers
whose effects were measured after the fact.

\begin{openingchart}{Estimated price changes after two completed U.S. consumer-product
mergers, by product: Pennzoil/Quaker State (motor oil) and Procter \& Gamble/Tambrands
(feminine hygiene), each measured against the prices of private-label products and of
other branded products over the same period (percent).}{Ashenfelter and Hosken (2008),
NBER Working Paper 13859, section 5}
\begin{tikzpicture}
\begin{axis}[omrecord, xbar, bar width=4pt, width=0.62\linewidth,
  height=0.42\linewidth, xmin=-2, xmax=10, ymin=-0.6, ymax=5.6,
  ytick={0,...,5}, yticklabels from table={figdata/bachelor-3/27-oligopoly-and-collusion/out/mergers.csv}{product},
  y tick label style={font=\scriptsize}, xlabel={Price change (\%)},
  legend pos=south east, legend style={font=\footnotesize}]
  \addplot[fill=omThm, draw=omThm, area legend] table[x=private, y=pos] {figdata/bachelor-3/27-oligopoly-and-collusion/out/mergers.csv};
  \addplot[fill=omDef!60, draw=omDef, area legend] table[x=branded, y=pos] {figdata/bachelor-3/27-oligopoly-and-collusion/out/mergers.csv};
  \legend{Private-label control, Branded control}
\end{axis}
\end{tikzpicture}
\end{openingchart}

% ledger: b3oc:ah.pqs
% ledger: b3oc:ah.pgt
% ledger: b3oc:ah
Against private-label products, every one of the six products rose in price by 4 to
8~\%; against other brands, the rises are smaller and one is negative. The two mergers
are among five that the authors chose because the agencies had found them most
problematic, so they bound the typical effect from above; four of the five raised
some prices, typically by 3 to 7~\%. The chart shows what a \omterm{def:b3:oligopoly-and-collusion:mergers}{merger retrospective}
measures and how much the answer depends on the control group.

\begin{recall}
Cournot and Bertrand competition and the Lerner index are the Year 1 volume's;
repeated games, the folk theorem and the collusion condition are the Year 2 volume's,
which also owns \emph{collusion} and \emph{tacit collusion}. The \omterm{def:b3:binary-outcomes:random-utility}{multinomial logit} is
\cref{ch:b3:binary-outcomes}'s, and difference in differences \cref{ch:b3:difference-in-differences}'s.
\end{recall}

\begin{notation}[Symbols of this chapter]\label{not:b3:oligopoly-and-collusion:symbols}
$s_j$ is product $j$'s market share and $s_0$ that of the outside option; $\delta_j$ its mean
utility; $\alpha$ the price sensitivity; $t$ a transport cost; $D_{12}$ the diversion ratio from
product 1 to product 2; $\mu$ a markup, price over marginal cost; $\beta$ the discount factor
of a repeated game (the Year 2 volume writes $\delta$, used here for mean utility).
\end{notation}

\section{Differentiated products}\label{sec:b3:oligopoly-and-collusion:differentiation}

\begin{definition}[Hotelling model; horizontal and vertical differentiation]\label{def:b3:oligopoly-and-collusion:hotelling}
\index{Hotelling model}\index{horizontal differentiation}\index{vertical differentiation}
In the \emph{Hotelling model}, or linear city, consumers are spread uniformly on a
segment, firms choose locations on it and then prices, and a consumer at distance $d$ from
a firm pays its price plus a transport cost. Products are \emph{horizontally
differentiated} when consumers at equal prices disagree on which they prefer, and
\emph{vertically differentiated} when all prefer the same one.
\end{definition}

\begin{proposition}[Hotelling with quadratic transport costs]\label{prop:b3:oligopoly-and-collusion:hotelling}
Let a unit mass of consumers lie on $[0, 1]$, two firms with marginal cost $c$ locate at $a$ and $1 - b$
($a, b \ge 0$, $a + b < 1$) and then set prices, and transport cost be $t d^2$. At the ends ($a = b = 0$) the
equilibrium prices are $c + t$; and each firm, choosing its location first, moves as far
from the other as it can.
\end{proposition}

\begin{proof}
The consumer indifferent between the firms is at
$\hat x = a + \tfrac{1}{2}(1 - a - b) + (p_2 - p_1)/(2t(1 - a - b))$. Firm 1 maximises $(p_1 - c)\hat x$; the first-order
conditions of the two firms solve to $p_1 = c + t(1 - a - b)(1 + (a - b)/3)$ and the symmetric $p_2$, which at
$a = b = 0$ give $c + t$. At these prices $\hat x = (3 + a - b)/6$ and firm 1's profit is
$\pi_1 = t(1 - a - b)(3 + a - b)^2/18$, so $\partial \pi_1/\partial a = t(3 + a - b)(-1 - 3a - b)/18 < 0$: moving towards the centre lowers profit,
because the price competition it provokes outweighs the demand it gains.
\end{proof}

\begin{omfigure}
\begin{tikzpicture}[econ, x=8cm, y=1cm]
  \draw[econ axes] (0,0) -- (1,0);
  \draw (0,0.08) -- (0,-0.08) node[below, font=\scriptsize] {0};
  \draw (1,0.08) -- (1,-0.08) node[below, font=\scriptsize] {1};
  \fill[omDef] (0,0) circle (2pt) node[above, font=\scriptsize] {firm 1, price $p_1$};
  \fill[omThm] (1,0) circle (2pt) node[above, font=\scriptsize] {firm 2, price $p_2$};
  \draw[curve, omDef] (0,0.5) parabola (0.62,1.65);
  \draw[curve, omThm] (1,0.75) parabola (0.38,1.9);
  \draw[guide] (0.54,0) -- (0.54,1.38);
  \node[below, font=\scriptsize] at (0.54,0) {$\hat x$};
  \node[font=\scriptsize, omDef] at (0.2,1.3) {$p_1 + t x^2$};
  \node[font=\scriptsize, omThm] at (0.82,1.4) {$p_2 + t(1 - x)^2$};
\end{tikzpicture}
\par\smallskip{\small The linear city. Each curve is the full price of buying from one
firm, for a consumer at $x$. Consumers left of $\hat x$ buy from firm 1, the rest from firm 2;
a lower $p_1$ shifts firm 1's curve down and moves $\hat x$ right.}
\end{omfigure}

\begin{proposition}[Bertrand with logit demand]\label{prop:b3:oligopoly-and-collusion:logit}
With logit demand $s_j = e^{\delta_j - \alpha p_j}/(1 + \sum_k e^{\delta_k - \alpha p_k})$, a firm selling product $j$ alone sets the
markup $p_j - c_j = 1/(\alpha(1 - s_j))$. A firm selling a set $J$ of products sets the same markup on
each, $1/(\alpha(1 - s_J))$, where $s_J$ is the share of the whole set.
\end{proposition}

\begin{proof}
Differentiating, $\partial s_j/\partial p_j = -\alpha s_j(1 - s_j)$ and $\partial s_j/\partial p_k = \alpha s_j s_k$ for $k \ne j$. For a single
product the first-order condition of $(p_j - c_j)s_j$ is $s_j - (p_j - c_j)\alpha s_j(1 - s_j) = 0$. For a set, the
condition for $k \in J$ is $s_k + \sum_{j \in J}(p_j - c_j)\,\partial s_j/\partial p_k = 0$. Try a common markup $m$: the sum is
$m\alpha s_k(-(1 - s_k) + \sum_{j \in J, j \ne k} s_j) = -m\alpha s_k(1 - s_J)$, so the condition holds at $m = 1/(\alpha(1 - s_J))$.
\end{proof}

Prices are strategic complements: a rival's higher price raises a firm's share and
so its markup. The proposition is the engine of a \omterm{def:b3:oligopoly-and-collusion:mergers}{merger simulation}. Owning both
products raises the share in the denominator from $s_j$ to $s_J$, and with it the markup; the
rivals then raise their prices too.

\begin{model}{Price competition with differentiated products}
\assumes Firms set prices simultaneously (Bertrand); demand comes from a
discrete-choice model; costs do not change with ownership.
\predicts Markups fall with the price sensitivity and with the complement of the
owner's share; a merger of close substitutes raises the prices of both and,
by less, of rivals.
\tested Simulations of five U.S. airline mergers of the 1980s, compared with the
prices observed afterwards, did not generally forecast accurately; most of the
unexplained change came from supply-side effects and from conduct other than the
assumed one. % ledger: b3oc:peters
\breaks When firms do not price as Bertrand competitors --- when they coordinate,
or set prices with an eye on the future --- or when the merger changes costs.
\end{model}

\section{Entry and market structure}\label{sec:b3:oligopoly-and-collusion:entry}

\begin{definition}[Entry threshold]\label{def:b3:oligopoly-and-collusion:entry-threshold}
\index{entry threshold}
The \emph{entry threshold} $S_N$ is the smallest market size, in customers, at which $N$
firms can each cover their fixed cost.
\end{definition}

\begin{proposition}[Thresholds per firm]\label{prop:b3:oligopoly-and-collusion:thresholds}
If each of $N$ firms earns variable profit $v_N$ per customer on a share $1/N$ of a market of size
$S$ and pays a fixed cost $F$, then $S_N = NF/v_N$. The threshold per firm, $S_N/N$, rises with $N$ exactly as long
as entry lowers the variable profit per customer; it stops rising once an extra firm
no longer changes conduct.
\end{proposition}

\begin{proof}
Each firm earns $(S/N)v_N - F$, which is zero at $S = NF/v_N$. Then $S_N/N = F/v_N$, which rises with $N$ if and
only if $v_N$ falls.
\end{proof}

\begin{model}{Entry and market structure}
\assumes Isolated markets differ in size; firms are alike and enter while profitable.
\predicts The number of firms grows with market size; the ratios $S_{N+1}/S_N$ reveal how fast
competition intensifies with $N$.
\tested In geographically isolated U.S. towns with one, two or a few doctors,
dentists, druggists, plumbers and tyre dealers, competitive conduct changed quickly as
the number of incumbents increased: the per-firm thresholds settle after the first
few entrants. % ledger: b3oc:br
\breaks When firms compete on quality or variety, so that an entrant changes the
products rather than the prices.
\end{model}

\section{Estimating market power}\label{sec:b3:oligopoly-and-collusion:markups}

Two routes lead to a markup. The demand route estimates demand, as in
\cref{ch:b3:binary-outcomes}, assumes a form of conduct and reads the markup off the
first-order conditions. The production route needs no demand and no conduct.

\begin{definition}[Production-approach markup estimator]\label{def:b3:oligopoly-and-collusion:markup-estimator}
\index{production-approach markup estimator}
The \emph{production-approach markup estimator} sets a firm's markup equal to the
output elasticity of a variable input divided by that input's share of revenue.
\end{definition}

\begin{proposition}[Markup from a variable input]\label{prop:b3:oligopoly-and-collusion:production}
If a firm minimises cost and can adjust input $V$ freely at price $P_V$, its markup $\mu = P/\lambda$, with $\lambda$
its marginal cost, equals $\theta_V/\alpha_V$, where $\theta_V = (\partial Q/\partial V)(V/Q)$ is the output elasticity of $V$ and
$\alpha_V = P_V V/(PQ)$ its share of revenue.
\end{proposition}

\begin{proof}
Cost minimisation sets $P_V = \lambda\,\partial Q/\partial V$, so $\theta_V = P_V V/(\lambda Q)$. Multiplying by $P/P$:
$\theta_V = (P/\lambda)\,P_V V/(PQ) = \mu\,\alpha_V$.
\end{proof}

The estimator needs an output elasticity, from an estimated production function, and
a decision about which costs are variable. That decision drives the result.

\begin{record}[Rising markups?]
With firm-level accounts for U.S. listed firms since 1950, and the cost of goods sold
as the variable input, average markups were stable until 1980 and then rose from 18~\%
above marginal cost to 67~\%, mostly within industries and through a sharp rise in
high-markup firms. % ledger: b3oc:de
With the same accounts and operating expenses as the variable input --- adding
marketing and management, a rising share of costs --- the rise since the 1980s is
modest and within historical variation, and calibrations that correct for listed
firms being unrepresentative give a flat or falling aggregate markup. % ledger: b3oc:traina
The two estimates use the same proposition. They differ in which costs they treat as
variable, which the accounts do not settle.
\end{record}

\section{Tacit collusion}\label{sec:b3:oligopoly-and-collusion:collusion}

\begin{proposition}[The collusion condition]\label{prop:b3:oligopoly-and-collusion:collusion}
In a repeated Bertrand market with $n$ identical firms and discount factor $\beta$, collusion at the
monopoly price, enforced by reverting forever to marginal-cost pricing after a
deviation, is sustainable if and only if $\beta \ge 1 - 1/n$.
\end{proposition}

\begin{proof}
Colluding gives $\pi^m/n$ each period, worth $\pi^m/(n(1 - \beta))$. Undercutting slightly takes the whole
monopoly profit $\pi^m$ once and zero after. Collusion holds iff $\pi^m/(n(1 - \beta)) \ge \pi^m$, that is $\beta \ge 1 - 1/n$.
\end{proof}

\begin{omchart}{The smallest discount factor at which collusion is sustainable in a
repeated Bertrand market, by the number of firms.}{The collusion condition of this
section}
\begin{tikzpicture}
\begin{axis}[omtime, width=0.6\linewidth, height=0.34\linewidth, xmin=1, xmax=21,
  ymin=0.4, ymax=1, xlabel={Number of firms $n$}, ylabel={$1 - 1/n$}]
  \addplot[omThm, mark=*, mark size=1.2pt] table[x=n, y=delta] {figdata/bachelor-3/27-oligopoly-and-collusion/out/collusion.csv};
\end{axis}
\end{tikzpicture}
\end{omchart}
% figdata: collusion

The condition extends in two directions that matter for detection. When demand is
temporarily high, the gain from deviating is larger relative to the future, so
collusive prices must fall in booms to hold the agreement together. When capacities
are unequal, the firm with the most spare capacity gains most from deviating, so
mergers that make capacities more equal can make collusion easier.

\begin{definition}[Cartel; leniency programme]\label{def:b3:oligopoly-and-collusion:cartel}
\index{cartel}\index{leniency programme}
A \emph{cartel} is an explicit agreement among competitors to fix prices, share
markets or rig bids. A \emph{leniency programme} grants immunity from, or a reduction
of, the penalty to the first cartel member that reports the cartel to the authority
and cooperates.
\end{definition}

\section{Detecting cartels}\label{sec:b3:oligopoly-and-collusion:detection}

\begin{measure}[A variance screen]
Collusion often shows in the movement of prices rather than their level. Around the
collapse of a bid-rigging conspiracy, the mean price fell by 16~\% while its standard
deviation rose by more than 200~\%. % ledger: b3oc:afgt
A screen searches data for pockets of unusually stable prices. Applied to retail
gasoline around Louisville from 1996 to 2002, it found none. A screen is a first
filter: stable prices can also come from stable costs, and a \omterm{def:b3:oligopoly-and-collusion:cartel}{cartel} that knows the
screen can add noise.
\end{measure}

\section{Mergers and what retrospectives found}\label{sec:b3:oligopoly-and-collusion:mergers}

\begin{definition}[Merger simulation; upward pricing pressure; merger retrospective]\label{def:b3:oligopoly-and-collusion:mergers}
\index{merger simulation}\index{upward pricing pressure}\index{merger retrospective}
A \emph{merger simulation} computes the post-merger equilibrium of an estimated
demand-and-conduct model with the new ownership. \emph{Upward pricing pressure} is
the net incentive of a merged firm to raise the price of one product, at pre-merger
prices. A \emph{merger retrospective} measures a completed merger's effect on prices
against a control.
\end{definition}

\begin{proposition}[Upward pricing pressure]\label{prop:b3:oligopoly-and-collusion:upp}
Let $D_{12} = -(\partial q_2/\partial p_1)/(\partial q_1/\partial p_1)$ be the share of the sales that product 1 loses when its price
rises which go to product 2. After a merger that cuts product 1's marginal cost by a share $E$, the merged firm wants
to raise $p_1$ at pre-merger prices if and only if $D_{12}(p_2 - c_2) > Ec_1$.
\end{proposition}

\begin{proof}
The merged firm's first-order condition for $p_1$ is $q_1 + (p_1 - c_1')\,\partial q_1/\partial p_1 + (p_2 - c_2)\,\partial q_2/\partial p_1 = 0$, with
$c_1' = (1 - E)c_1$. Dividing by $-\partial q_1/\partial p_1$, it reads as the pre-merger condition with marginal cost
$c_1' + D_{12}(p_2 - c_2)$: each unit lost by product 1 now costs the margin on the $D_{12}$ units it sends to
product 2. That effective cost exceeds $c_1$, and the pre-merger price is too low, iff
$D_{12}(p_2 - c_2) > Ec_1$.
\end{proof}

The test needs only a diversion ratio and a margin, not a full model; it is a
first-order screen, holding the rivals' prices fixed and ignoring how much of a cost
change the firm passes on. A full simulation, as in the weekend problem, gives the
equilibrium.

% ledger: b3oc:ahw
Retrospectives of mergers in concentrated markets have repeatedly found price rises.
The record supports the model's direction: mergers in oligopolistic markets can raise
consumer prices. The magnitudes depend on the control group, as the opening chart
shows, and simulations made before the merger have missed them in both directions.

\section{Exercises}\label{sec:b3:oligopoly-and-collusion:exercises}

\begin{exercise}[$\star$]\label{exo:b3:oligopoly-and-collusion:1}
In the opening chart, find the median price change against each control group, and
the number of negative estimates.
\end{exercise}

\begin{exercise}[$\star$]\label{exo:b3:oligopoly-and-collusion:2}
In the \omterm{def:b3:oligopoly-and-collusion:hotelling}{Hotelling model} with firms at the ends, $c = 10$ and $t = 4$, find the prices and each firm's
profit.
\end{exercise}

\begin{exercise}[$\star$]\label{exo:b3:oligopoly-and-collusion:3}
With logit demand, $\alpha = 0.5$ per dollar and a share of 20~\%, what is a single-product firm's markup?
\end{exercise}

\begin{exercise}[$\star$]\label{exo:b3:oligopoly-and-collusion:4}
What pattern of prices over time would contradict tacit collusion in a market?
\end{exercise}

\begin{exercise}[$\star\star$]\label{exo:b3:oligopoly-and-collusion:5}
Find the collusion threshold for 2 and for 5 firms. By how much may the interest rate
per period rise, from 0, before collusion among 5 firms breaks down?
\end{exercise}

\begin{exercise}[$\star\star$]\label{exo:b3:oligopoly-and-collusion:6}
Three symmetric logit firms with $\alpha = 1$ each have a share of 0.25. Two merge. Holding shares fixed
at 0.25, what markup does the merged firm want? Why does the true equilibrium markup
differ?
\end{exercise}

\begin{exercise}[$\star\star$]\label{exo:b3:oligopoly-and-collusion:7}
The diversion ratio from product 1 to product 2 is 0.3, product 2's margin is 4 dollars and
product 1's marginal cost 10. Compute the \omterm{def:b3:oligopoly-and-collusion:mergers}{upward pricing pressure} on product 1, and the
cost saving that offsets it.
\end{exercise}

\begin{exercise}[$\star\star$]\label{exo:b3:oligopoly-and-collusion:8}
An invented table of 400 isolated towns gives population and the number of dentists.
Estimate the \omterm{def:b3:oligopoly-and-collusion:entry-threshold}{entry thresholds} and the thresholds per dentist. After how many dentists
does entry stop changing conduct?
\end{exercise}

\begin{exercise}[$\star\star$]\label{exo:b3:oligopoly-and-collusion:9}
What evidence would show that a measured rise in markups is an artefact of how costs
are classified?
\end{exercise}

\begin{exercise}[$\star\star\star$]\label{exo:b3:oligopoly-and-collusion:10}
Prove the logit markup formula for a single product and for a set of products.
\end{exercise}

\begin{exercise}[$\star\star\star$]\label{exo:b3:oligopoly-and-collusion:11}
Four products have logit demand with $\alpha = 1$, cost 1 and mean utilities 3.0, 2.8, 2.0 and 1.0. Simulate
the merger of products 1 and 2, and of products 1 and 4. Compare the price rises of product
1 with the \omterm{def:b3:oligopoly-and-collusion:mergers}{upward pricing pressure}.
\end{exercise}

\begin{exercise}[$\star\star\star$]\label{exo:b3:oligopoly-and-collusion:12}
Two \omterm{def:b3:oligopoly-and-collusion:cartel}{cartel} members each choose to stay silent or report. If neither reports, each
pays a fine $F$ with probability $q$; a sole reporter pays nothing and the other pays $F$; if both report,
each pays $F/2$. Show that reporting is a dominant strategy for any $q > 0$, and explain why a
\omterm{def:b3:oligopoly-and-collusion:cartel}{leniency programme} destabilises \omterm{def:b3:oligopoly-and-collusion:cartel}{cartels}.
\end{exercise}

\section{Problem: Three Bakeries in Brenmoor}\label{sec:b3:oligopoly-and-collusion:problem}

\begin{problem}[{Weekend problem --- three bakeries in Brenmoor}]\label{pb:b3:oligopoly-and-collusion:bakeries}
Brenmoor's 3{,}000 loaves a day are sold by three bakeries, or not bought at all. Demand
is logit with $\alpha = 0.8$ per dalen and mean utilities 2.5, 2.5 and 2.1; each bakery's marginal cost is
2 dalens a loaf. Bakeries 1 and 2 propose to merge. % figdata: bakeries_stats

\medskip\noindent\textbf{Part I --- the market before the merger.}
\begin{enumerate}
\item Write each bakery's share as a function of the three prices.
\item State the markup each bakery sets.
\item Find the equilibrium prices (iterate the markup formula).
\item Find the shares and the share of residents who buy no loaf.
\item Check the markup formula for bakery 1.
\item Find bakery 1's daily profit.
\item Find consumer surplus, $\ln(1 + \sum_j e^{\delta_j - \alpha p_j})/\alpha$ per resident, for the market.
\end{enumerate}

\medskip\noindent\textbf{Part II --- simulating the merger.}
\begin{enumerate}[resume]
\item Why does the merged firm raise both its prices?
\item Find the post-merger prices.
\item Find the shares.
\item By how much do the merged bakeries' prices rise, and bakery 3's?
\item Find the loss of consumer surplus per day.
\item Compare the merged firm's profit with the sum of the two bakeries' profits before.
\item Why does bakery 3 raise its price, although nothing changed for it directly?
\end{enumerate}

\medskip\noindent\textbf{Part III --- efficiencies.}
\begin{enumerate}[resume]
\item Find the diversion ratio from bakery 1 to bakery 2 at pre-merger prices.
\item Find bakery 2's margin.
\item Compute the \omterm{def:b3:oligopoly-and-collusion:mergers}{upward pricing pressure} on bakery 1.
\item What cost saving, as a share of marginal cost, would offset it by this test?
\item The full simulation says the merged bakeries' prices stay at the pre-merger level
only if their marginal cost falls to 1.27 dalens. What saving is that?
\item Why do the two answers differ?
\end{enumerate}

\medskip\noindent\textbf{Part IV --- a retrospective.} The merger goes ahead. Two years later
the share-weighted mean price of a loaf in Brenmoor is 0.40 dalens above its pre-merger
level; in a comparison town it has gone from 3.90 to 4.05.
\begin{enumerate}[resume]
\item Compute the difference-in-differences estimate.
\item Compute the simulation's prediction for the share-weighted mean price.
\item Compare the two.
\item Name two threats to the comparison.
\item What could explain a measured effect below the predicted one?
\item State the result: the merger's predicted and measured price effects, and the cost
saving that would have made it neutral.
\end{enumerate}
\end{problem}
